\documentclass[a4paper]{article}
\usepackage{amsfonts,amsmath,amssymb,graphicx}

\begin{document}

\noindent \large Auteur: Michael Zielinski \\
\noindent \large Studierichting: Informatica\\
\noindent \large Oefeningenreeks 3D nr. 7.2 \\

\normalsize

Bepaal de extrema van $f(x)=e^x\sin x$. \\

\section*{1. Eerste afgeleide}

$f'(x)=\dfrac{d}{dx}(e^x\sin x)$ \\

$=\left(\dfrac{d}{dx}e^x\right)\sin x
+ e^x\left(\dfrac{d}{dx}\sin x\right)$ \\

$=e^x\sin x + e^x\cos x$ \\

$=e^x\left(\sin x+\cos x\right)$ \\

\section*{2. Kritische punten}

We zoeken $f'(x)=0$. \\

$e^x\left(\sin x+\cos x\right)=0$ \\

Omdat $e^x>0$ voor alle $x$ geldt: \\

$\sin x+\cos x=0$ \\

$\Leftrightarrow\sin x=-\cos x$ \\

$\Leftrightarrow\tan x=-1$ \\

$\Leftrightarrow x_k=-\dfrac{\pi}{4}+k\pi,\;k\in\mathbb{Z}$ \\

\section*{3. Tweede afgeleide}

$f''(x)=\dfrac{d}{dx}\left(e^x(\sin x+\cos x)\right)$ \\

$=\left(\dfrac{d}{dx}e^x\right)(\sin x+\cos x)
+ e^x\left(\dfrac{d}{dx}(\sin x+\cos x)\right)$ \\

$=e^x(\sin x+\cos x) + e^x(\cos x-\sin x)$ \\

$=2e^x\cos x$ \\

Daarom: \\

$f''(x_k)=2e^{x_k}\cos x_k
=2e^{x_k}\dfrac{\sqrt{2}}{2}(-1)^k
=\sqrt{2}e^{x_k}(-1)^k$ \\

\section*{4. Classificatie en specificatie}

Uit $x_k=-\dfrac{\pi}{4}+k\pi$ onderscheiden we:

\begin{itemize}
\item Stel $k=2n$ (even). \\

  $x_{2n}=-\dfrac{\pi}{4}+2n\pi$. \\

  Dan $(-1)^{2n}=+1$, dus $f''(x_{2n})>0$. \\

  $\Rightarrow$ lokaal minimum bij \\

  $x=-\dfrac{\pi}{4}+2n\pi,\;n\in\mathbb{Z}$ \\

\item Stel $k=2n+1$ (oneven). \\

  $x_{2n+1}=-\dfrac{\pi}{4}+(2n+1)\pi
  =-\dfrac{\pi}{4}+2n\pi+\pi
  =\dfrac{3\pi}{4}+2n\pi$. \\

  Dan $(-1)^{2n+1}=-1$, dus $f''(x_{2n+1})<0$. \\

  $\Rightarrow$ lokaal maximum bij \\

  $x=\dfrac{3\pi}{4}+2n\pi,\;n\in\mathbb{Z}$ \\

\end{itemize}

\section*{5. Waarden in de extrema}

$f(x_k)=e^{x_k}\sin x_k
=e^{x_k}\left(-\dfrac{\sqrt{2}}{2}(-1)^k\right)
=-\dfrac{\sqrt{2}}{2}(-1)^k e^{x_k}$ \\

\begin{itemize}
\item Lokaal minimum ($k=2n$): \\

  $x=-\dfrac{\pi}{4}+2n\pi,\;f(x)=-\dfrac{\sqrt{2}}{2}\,e^{-\tfrac{\pi}{4}+2n\pi}$. \\

\item Lokaal maximum ($k=2n+1$): \\

  $x=\dfrac{3\pi}{4}+2n\pi,\;f(x)=+\dfrac{\sqrt{2}}{2}\,e^{\tfrac{3\pi}{4}+2n\pi}$. \\

\end{itemize}

\begin{thebibliography}{9}
% \bibitem{ref} Referentie
\end{thebibliography}

\end{document}
